\documentclass[12pt, letterpaper]{article}

%-------Packages---------
\usepackage{amssymb,amsfonts}
\usepackage{mathptmx}
\usepackage{amsthm}
\usepackage{graphicx}
\usepackage[all,arc]{xy}
\usepackage[colorlinks=true, allcolors=blue]{hyperref}
\usepackage{enumerate}
\usepackage{bbm}
\usepackage{dsfont}
\usepackage{mathrsfs}
\usepackage{tikz-cd}
\usepackage{setspace}
\usepackage{import}
\usepackage[utf8]{inputenc}
\usepackage{hyperref}
\usepackage{tikz}
\usepackage{float}
\usepackage{mdframed}
\usepackage{fancyhdr}
\usepackage[nointegrals]{wasysym}

\usepackage[left=1in,top=1in,right=1in,bottom=1in]{geometry}

\usepackage{mathtools}
\usepackage{setspace}
\usepackage{cleveref}
\usepackage{multicol}

%--------Theorem Environments--------
%theoremstyle{plain} --- default
\newmdtheoremenv{theorem}{Theorem}[subsection]
\newmdtheoremenv{corollary}[theorem]{Corollary}
\newtheorem{proposition}[theorem]{Proposition}
\newmdtheoremenv{lemma}[theorem]{Lemma}
\newtheorem{conj}[theorem]{Conjecture}
\newtheorem{quest}[theorem]{Question}

\theoremstyle{definition}
\newmdtheoremenv{definition}[theorem]{Definition}
\newmdtheoremenv{definitions}[theorem]{Definitions}
\newtheorem{example}[theorem]{Example}
\newtheorem{algorithm}[theorem]{Algorithm}
\newtheorem{rmk}[theorem]{Remark}
\newtheorem{exercise}[theorem]{Exercise}

%----------- my commands -------------
\newcommand{\C}{\mathbb{C}}
\newcommand{\R}{\mathbb{R}}
\newcommand{\Q}{\mathbb{Q}}
\newcommand{\Z}{\mathbb{Z}}
\newcommand{\N}{\mathbb{N}}
\renewcommand{\P}{\mathbb{P}}
\newcommand{\contr}{\Rightarrow \Leftarrow}
\newcommand{\HRule}[1]{\rule{\linewidth}{#1}}
\newcommand{\s}{\(\,\)}
\renewcommand{\t}[1]{\text{#1}}
\newcommand{\ang}[1]{\left\langle #1 \right\rangle}

% Meta data
\setstretch{1.2}

\begin{document}

\pagestyle{fancy}
\fancyhf{}
\fancyhead[LOE]{\slshape{\textbf{Vincent Ferrara}}}
\fancyhead[ROE]{\thepage}

\subsection*{Problem 1}
Let $S_n=X_1+\dots+X_n$
\begin{align*}
    \varphi_{S_n/n}(t)=E(e^{itS_n/n})&=\prod_{k=1}^{n}E(e^{itX_k/n})\\
    &=\prod_{k=1}^{n}\varphi_{X_k}(t/n)\\
    &=\prod_{k=1}^{n}(1+\dfrac{it \mu}{n}+O(\dfrac{\t{Var}(X_k)t^2}{n^2}))\\
    &=\prod_{k=1}^{n}(1+O(\dfrac{\t{Var}(X_k)t^2}{n^2}))
    \intertext{Then take limit as $n \rightarrow \infty$}
    \lim_{n\rightarrow \infty}\varphi_{S_n/n}(t)=\lim_{n\rightarrow \infty}\prod_{k=1}^{n}(1+O(\dfrac{\t{Var}(X_k)t^2}{n^2}))=1=e^0
    \intertext{Since $O(\dfrac{\t{Var}(X_k)t^2}{n^2})$ goes to zero}
\end{align*}
Thus $S_n/n \stackrel{(d)}{\longrightarrow} 0$

\newpage
\subsection*{Problem 2}
Shown in class $\varphi_X(t)=e^{-t^2/2}$ and $f_X(t)=\dfrac{1}{\sqrt{2\pi}}e^{-x^2/2}$\\
Thus $\varphi_{XY}(t)=E(e^{itXY})=E(E_Y(e^{itXY}|X))$ thus since they are independent $E_Y(e^{itXY}|X)=\varphi_Y(tX)$\\
$E(\varphi_Y(Xt))=E(e^{-t^2X^2/2})$
\begin{align*}
    E(e^{-t^2X^2/2})&=\int_{-\infty}^{\infty}e^{-t^2x^2/2}f_X(x)dx\\
    &=\dfrac{1}{\sqrt{2\pi}}\int_{-\infty}^{\infty}e^{-t^2x^2/2}e^{-x^2/2}dx\\
    &=\dfrac{1}{\sqrt{2\pi}}\int_{-\infty}^{\infty}e^{-1/2x^2(1+t^2)}dx\\
    &=\dfrac{1}{\sqrt{2\pi}}\sqrt{\dfrac{2\pi}{1+t^2}}\\
    &=\dfrac{1}{\sqrt{1+t^2}}
\end{align*}
\phantom{text}\\
Thus for the next part since they are all independent we can just use the function we found in the first part then multiply them together thus we get $\varphi(t)=\dfrac{1}{1+t^2}$\\
\begin{align*}
    \varphi_X(t)=E(e^{itX})&=\int_{-\infty}^{\infty}e^{itx}(1/2)e^{-|x|}dx\\
    &=(1/2)\int_{-\infty}^{0}e^{itx}e^{x}dx+(1/2)\int_{0}^{\infty}e^{itx}e^{-x}dx\\
    &=\dfrac{1}{2(it-1)}+\dfrac{1}{2(1+it)}\\
    &=\dfrac{1}{1+t^2}
\end{align*}
Thus since the characteristic of $f(x)=(1/2)e^{-|x|}$ is the same we know that the pdf of $X_1X_2+X_3X_4$ is $(1/2)e^{-|x|}$

\newpage
\subsection*{Problem 3}
\begin{align*}
    \varphi_Z(t)=E(e^{itZ})&=\int_{-\infty}^{\infty}\int_{-\infty}^{\infty}e^{itz}f(x,z+2x)dxdz\\
    &=(1/2\pi)\int_{-\infty}^{\infty}\int_{-\infty}^{\infty}e^{itz}exp\left(  -\dfrac{5x^2-4xz-8x^2+z^2+4xz+4x^2}{2}\right)dxdz\\
    &=(1/2\pi)\int_{-\infty}^{\infty}\int_{-\infty}^{\infty}e^{itz}e^{-(1/2)(x^2+z^2)}dxdz\\
    &=(1/2\pi)\int_{-\infty}^{\infty}e^{itz-(1/2)z^2}\int_{-\infty}^{\infty}e^{-(1/2)(x^2)}dxdz\\
    &=(1/2\pi)\dfrac{\sqrt{\pi}}{\sqrt{1/2}}\int_{-\infty}^{\infty}e^{-(-itz+(1/2)z^2)}dz\\
    &=(1/2\pi)\dfrac{\sqrt{\pi}}{\sqrt{1/2}}\dfrac{\sqrt{\pi}}{\sqrt{1/2}}e^{-t^2/2}\\
    &=e^{-(1/2)t^2}
\end{align*}
Thus we know that $Z$ is an N(0,1) since it has the same character.
\end{document}